\chapter{Backup, recovery, migration, and replication}
\label{ops:recovery}

\section{Define what must be recoverable}

A backup is useful only if it can be restored with the available executable,
keys, storage capacity, and operational instructions. Decide which data loss
is acceptable and how long recovery may take. A periodic base backup protects
a point in time. A WAL archive can extend recovery beyond that base. An
asynchronous standby can shorten some failover steps but may lag committed
primary work. These mechanisms have different failure boundaries.

Keep the source revision/build, configuration, external keys, and restore
instructions alongside the backup inventory, with keys protected separately
from the captured database artifacts. A copy on the same failed disk does
not create an independent recovery location. Test a restore into a fresh
directory instead of assuming that a successful copy command proves recovery.

\begin{figure}[htbp]
\centering
\begin{tikzpicture}[node distance=12mm and 11mm]
\node[box,text width=37mm] (init) {Create database\\\texttt{--init}};
\node[strong,text width=37mm,right=of init] (serve) {Open server\\Connect clients};
\node[box,text width=37mm,right=of serve] (stop) {Graceful shutdown\\Release file locks};
\node[box,text width=37mm,below=of init] (restore) {Restore to new path\\Validate separately};
\node[box,text width=37mm,below=of serve] (backup) {Offline backup\\Keep matching keys};
\node[box,text width=37mm,below=of stop] (check) {Offline check\\or page-size rewrite};
\draw[arrow] (init) -- (serve);\draw[arrow] (serve) -- (stop);\draw[arrow] (stop) -- (check);
\draw[arrow] (stop.south) -- ++(0,-4mm) -| (backup.north);
\draw[arrow] (backup) -- (restore);\draw[arrow] (restore.west) -- ++(-5mm,0) |- (init.west);
\draw[arrow] (check.east) -- ++(5mm,0) |- (serve.north);
\end{tikzpicture}
\caption{A safe teaching workflow. A restore is a new candidate database that must be checked before it replaces the working copy. Advanced replication has additional base, WAL and fencing requirements.}
\label{fig:database-lifecycle}
\end{figure}


\section{Create and verify a base backup}
\index{backup}

\subsection{The command-line backup is an offline operation}

The base-backup command opens the database directly and takes its normal
database ownership lock. Stop the server before invoking it. The internal
\code{runOpen} API can work with an already-open managed database, but that
does not make the public standalone CLI a network online-backup service.

Choose a destination that does not exist, including its staging counterpart:

\begin{lstlisting}
minisql-backup.exe backup C:\MiniSQL-Data\db_ACTUAL_ID `
  E:\MiniSQL-Backups\learning-base-001
\end{lstlisting}

\begin{lstlisting}[language=bash]
minisql-backup backup /srv/minisql/db_ACTUAL_ID \
  /backup/minisql/learning-base-001
\end{lstlisting}

MiniSQL captures authoritative metadata, catalog/security state, schema state,
optional statistics/encryption metadata, audit artifacts, WAL/checkpoint
state, tables, and indexes. It writes a directory containing captured files
and a checksummed \path{backup.manifest}. The report gives file count and
total bytes. The destination is not a single ZIP file and should be treated
as one complete backup unit.

Derived \path{*.heap-pages} scan directories are not needed for the backup;
restored databases rebuild them from authoritative table pages. Runtime
configuration, executable binaries, external raw keys, TLS private keys, and
ordinary server log directories are not all automatically included in this
database capture. Inventory those separately.

\subsection{Capacity and memory limitations}

The current capture code reads file images into memory and retains the
captured file collection while constructing the backup. It is not the
streaming algorithm used by the checker. A large database can therefore need
substantial working memory in addition to destination disk space. Encrypted
backup temporarily retains captured and encrypted representations as well.
Measure the largest real backup, not just a tiny demonstration database.

The explicit per-file safety limit is 4,294,967,295 bytes. The manifest uses
wide lengths, but the current implementation still applies this U32-sized
file cap. Its total safety bound is
1,152,921,504,606,846,975 bytes and its entry-count cap is U32; these format
bounds are not claims that the in-memory implementation can practically back
up that amount of data. Relative manifest paths are limited to 240 bytes and
validated against traversal/absolute-path forms.

\sourcecode{src/minisql/tools/backup.ml}{captureDatabase, runOpen, manifest limits, and path validation.}

\section{Restore a base into a new directory}
\index{restore}

\begin{lstlisting}
minisql-backup.exe restore E:\MiniSQL-Backups\learning-base-001 `
  C:\MiniSQL-Restores\learning-restore-001
minisql-check.exe C:\MiniSQL-Restores\learning-restore-001
\end{lstlisting}

The target and its \path{.restore-new} staging path must not exist. Restore
validates captured file lengths/checksums, stages the database, opens it for
recovery/verification, checks its database identity, and publishes the
validated target. It does not overwrite an existing live database in place.
The restored database retains the backup's logical database identity; the
chosen directory name is an operational path, not a new SQL database name.

After structural checking, start the restored copy on a separate unused port
and run semantic application checks: representative row counts, important
records, totals computed with the currently documented type semantics, and
relationship checks. Do not expose both the old and restored writable copy to
the same application by accident. Choose a single authoritative copy and
update application connection configuration deliberately.

If a restore fails, preserve the diagnostic and the original backup. A
leftover staging directory can be evidence of the failure. Inspect it before
removing it to retry, and never mistake a staging directory for a fully
published restore. Existing-target refusal is a deliberate guard against
overwriting prior data or a previous recovery attempt.

\sourcecode{src/minisql/tools/backup.ml}{restore and post-copy identity/consistency validation.}

\subsection{Automate the restore drill}

The repository's recovery drill performs backup, restore, and checker stages,
records timings and bounded command output, and writes a JSON report:

\begin{lstlisting}[language=bash]
python3 tools/recovery/production_recovery_drill.py \
  /srv/minisql/db_ACTUAL_ID \
  /backup/minisql/drill-base-001 \
  /srv/minisql-restores/drill-001 \
  --bin-dir build/bin-linux \
  --report build/recovery/drill-001.json
\end{lstlisting}

The default binary directory is \path{build/bin}; Linux users must select
\path{build/bin-linux} unless they have installed the programs elsewhere.
The default report path is \path{build/recovery/recovery-drill.json}.
Both output destinations must be new. This runner does not start and stop
your live server for you; arrange its offline prerequisite first.

\sourcecode{tools/recovery/production_recovery_drill.py}{Arguments, three stages, and evidence report.}

\section{Encrypt a backup independently of live TDE}

Generate a distinct 32-byte backup key outside the backup tree and create an
encrypted backup:

\begin{lstlisting}
minisqld.exe --generate-encryption-key D:\MiniSQL-Keys\backup-001.key
minisql-backup.exe backup-encrypted C:\MiniSQL-Data\db_ACTUAL_ID `
  E:\MiniSQL-Backups\learning-encrypted-001 `
  D:\MiniSQL-Keys\backup-001.key
minisql-backup.exe restore-encrypted `
  E:\MiniSQL-Backups\learning-encrypted-001 `
  C:\MiniSQL-Restores\learning-encrypted-001 `
  D:\MiniSQL-Keys\backup-001.key
\end{lstlisting}

Each captured file is encrypted/authenticated independently and remains
covered by the manifest. Wrong keys or modified encrypted payloads fail
closed. If live TDE is enabled, the backup process rewraps its DEK for the
backup provider. On restore, the same key bytes can be supplied from a new
machine-local path; MiniSQL authenticates the exported envelope and rewraps
the DEK to that restore provider identity before opening the database.

Keep the backup key available for the retained backup's lifetime. Rotating the
live KEK does not magically update already-created backups. The normal
\code{restore} command is not a substitute for \code{restore-encrypted},
and a password reset is not an encrypted-backup recovery method. The backup
key is a binary file argument, not a SQL user's password.

\sourcecode{src/minisql/tools/backup.ml}{Encrypted record format, runOpenEncrypted, and restoreEncrypted.}

\section{WAL archives and point-in-time recovery}
\index{PITR}\index{WAL archive}\index{recovery}
\begin{figure}[htbp]\centering
\begin{tikzpicture}[node distance=10mm and 12mm]
\node[strong,text width=47mm] (base) {Verified base backup\\Schema and physical identity};
\node[box,text width=47mm,right=of base] (wal) {Contiguous archived WAL\\Same base and schema contract};
\node[strong,text width=47mm,below=of base] (restore) {Restore to a new directory\\Replay to chosen LSN / latest};
\node[box,text width=47mm,right=of restore] (key) {Matching external keys\\when encryption is used};
\node[box,text width=47mm,below=of restore] (verify) {Offline consistency check\\Then application-level checks};
\node[strong,text width=47mm,right=of verify] (publish) {Select authoritative copy\\Update client endpoint};
\draw[arrow] (base) -- (restore);\draw[arrow] (wal.south) -- ++(0,-4mm) -| (restore.north);
\draw[arrow] (key) -- (restore);\draw[arrow] (restore) -- (verify);\draw[arrow] (verify) -- (publish);
\end{tikzpicture}
\caption{The evidence needed for point-in-time recovery. A missing WAL range, incompatible schema base or unavailable key cannot be repaired by choosing a later target.}
\label{fig:recovery-chain}
\end{figure}


\subsection{Understand the archive model}

A WAL archive contains a verified base backup and successive complete durable
WAL prefixes. Its manifest records database identity, page size, generation,
base-end LSN, latest-end LSN, WAL file name, length, and checksum. LSN is a
byte-position boundary in this WAL model, not a timestamp or SQL transaction
number.

Initialize an archive while the primary is stopped:

\begin{lstlisting}
minisql-backup.exe archive-init C:\MiniSQL-Data\db_ACTUAL_ID `
  E:\MiniSQL-Archives\learning-001
minisql-backup.exe archive-verify E:\MiniSQL-Archives\learning-001
\end{lstlisting}

Start the primary and export a durable prefix while it remains online:

\begin{lstlisting}
minisql-backup.exe archive-wal-live C:\MiniSQL-Data\db_ACTUAL_ID `
  E:\MiniSQL-Archives\learning-001
\end{lstlisting}

The live form reads the durable WAL marker and validates that the captured
prefix ends at a complete record boundary. It verifies continuity against
the previous archive prefix. If no new bytes are durable, it can return the
same generation. The \code{archive-wal} form opens the database normally
and is an offline form. Use only one archive writer for a given archive
directory.

Every generation stores a complete prefix, not only the delta from the
previous generation. Older files are retained so a verifier that already
read an older manifest can still find its referenced file. Archive growth
can therefore be substantial. Create a fresh verified base periodically and
retire old archive sets under a retention policy only after their recovery
requirements have expired.

\warning{The archive follows a fixed base schema}{The current shipping model
is intended for DML on the already archived schema/security state. After DDL,
DCL, VACUUM, page-size migration, key/envelope transitions, or operations that
rewrite or rewind physical state, stop the replication workflow and create a
new verified base. A lost WAL-prefix relationship is diagnosed; editing the
manifest does not repair it.}

\sourcecode{src/minisql/tools/backup.ml}{archiveInit, snapshotDurableWalLive, archiveWalLive, and verifyArchive.}
\sourcecode{docs/formats/wal-archive-v1.md}{Archive and standby metadata fields.}

\subsection{Automatic checkpoints can break archive continuity}

The current automatic checkpoint implementation resets the physical WAL to
zero when it reaches \code{runtime.checkpointWalBytes}, 64 MiB by default.
It has no archive-aware retention handshake. A subsequent live export then
finds that the new WAL is not an extension of the archived prefix and rejects
the export. Normal DML can reach this boundary even when no operator has
issued VACUUM or changed the schema.

Do not assume that either sidecar automatically rebases across this reset.
The hot-primary exporter propagates the failure. The HA controller's
replication error also propagates, and its cleanup stops the managed primary.
For a qualified archival workflow, plan a new base before a reset, monitor
WAL growth/checkpoint resets, and select a threshold consistent with that
maintenance window where the chosen launcher exposes configuration. The
current fenced controller uses the positional server mode and does not offer
a checkpoint-threshold argument. Increasing a threshold postpones the issue;
it does not implement archive-aware WAL retention, and the archive file-size
and memory limits still apply.

\sourcecode{src/minisql/server/database_manager.ml}{checkpointIfNeeded rewinds WAL without an archive-retention handshake.}
\sourcecode{tools/replication/minisql_ha_controller.py}{Replication failure propagation and process cleanup.}

\subsection{Restore to latest or an exact record boundary}

\begin{lstlisting}
minisql-backup.exe restore-pitr E:\MiniSQL-Archives\learning-001 `
  C:\MiniSQL-Restores\learning-latest latest
minisql-check.exe C:\MiniSQL-Restores\learning-latest
\end{lstlisting}

For an earlier recovery target, replace \code{latest} with an exact numeric
LSN taken from valid archive/WAL evidence. The target must lie between the
base-end and latest-end LSN and must be a complete record boundary. A number
in the middle of a record is rejected. There is no public
\code{restore-until-timestamp} switch; mapping a business event to an LSN
requires operational evidence outside this command.

Recovery replays committed transactions within the selected prefix. A
transaction whose commit has not been included is not made committed merely
because some of its records precede the target. When post-base authoritative
table pages have been replayed, the base indexes are marked dirty and rebuilt
before the restored generation is published. An exact base-boundary restore
can retain the already verified base indexes.

The restore destination and its \path{.pitr-stage} must be new. Preserve the
original primary if investigating an accidental change; restore separately,
check the selected state, and only then decide how to resume application
traffic. PITR is a database-wide recovery operation, not a built-in single-row
undo command.

\sourcecode{src/minisql/tools/backup.ml}{walPrefixAt, restoreToLsn, restoreLatest, and repairRecoveredIndexes.}

\section{Materialize and refresh a read-only standby}
\index{standby}

\begin{lstlisting}
minisql-backup.exe standby-materialize `
  E:\MiniSQL-Archives\learning-001 C:\MiniSQL-Standby\slot-a
minisqld.exe --serve-standby C:\MiniSQL-Standby\slot-a 7433 16
\end{lstlisting}

Materialization restores the latest verified archive state and writes a
standby marker before publishing the target. The marker records database
identity, archive generation, and applied LSN. The ordinary writable open path
rejects an unpromoted standby. The standby server permits reads and rejects
attempts to use it as a writable primary.

This describes ordinary mutation dispatch, not a complete safety guarantee for
all expression paths. In this revision, NEXTVAL can be materialized during
streaming-plan probing or EXPLAIN outside the ordinary read-only checks.
Exclude sequence-consuming expressions from standby workloads; do not use
this known gap to advance a standby's state. See the transaction and sequence
chapters for the observed behavior and restrictions.

To advance a manually managed standby, stop its server first, export the
latest primary WAL, and refresh:

\begin{lstlisting}
minisql-backup.exe standby-refresh E:\MiniSQL-Archives\learning-001 `
  C:\MiniSQL-Standby\slot-a
\end{lstlisting}

\warning{Stop the standby before refresh}{The direct refresh routine replaces
WAL/checkpoint files before opening the standby through the database manager.
Do not rely on a later lock conflict to make an attempted live refresh safe.
The public command is for a stopped standby. Use the alternating-slot sidecar
for continued read service while a separate generation is prepared.}

Refresh validates the archive identity and monotonic generation/LSN, checks
the prefix relationship, replays the new state, repairs derived indexes, and
updates the standby marker. A standby can be behind the primary even while
its server answers queries normally. Monitor applied LSN and source LSN;
availability of the read endpoint is not proof that it is current.

\sourcecode{src/minisql/tools/backup.ml}{materializeStandby, refreshStandby, and standby markers.}
\sourcecode{src/minisql/server/database_manager.ml}{Writable-open rejection of unpromoted standbys.}

\section{Continuous shipping with the hot-replica sidecar}
\index{replication}

The Python sidecar automates durable-prefix export and maintains two standby
slots. A loopback TCP switch sends new connections to the current verified
generation while another slot is prepared. Existing connections remain
associated with their selected backend; they are not silently transferred
mid-session to a different database generation. Backend retirement waits for
clients according to the sidecar's lifecycle/timeout policy.

After creating the initial archive, run the primary exporter in a separate
process:

\begin{lstlisting}[language=bash]
python3 tools/replication/minisql_hot_replica.py primary \
  --database /srv/minisql/db_ACTUAL_ID \
  --archive /srv/minisql-ha/archive \
  --backup-exe build/bin-linux/minisql-backup \
  --status-file /srv/minisql-ha/primary-status.json
\end{lstlisting}

Run the standby sidecar against that archive:

\begin{lstlisting}[language=bash]
python3 tools/replication/minisql_hot_replica.py standby \
  --archive /srv/minisql-ha/archive \
  --slot-root /srv/minisql-ha/slots \
  --backup-exe build/bin-linux/minisql-backup \
  --server-exe build/bin-linux/minisqld \
  --listen-port 7433 \
  --status-file /srv/minisql-ha/standby-status.json
\end{lstlisting}

The reference proxy binds only to loopback because it fronts a trusted
standby listener. This script does not distribute archives to another machine,
provision network storage, or add remote authentication by itself. Arrange
the intended shared/transferred archive and its access controls outside the
script. Only one process may write a particular archive.

\begin{longtable}{p{.34\linewidth}p{.59\linewidth}}
\toprule
Sidecar option & Meaning / default \\
\midrule
\endhead
Both: \code{--archive}, \code{--backup-exe} & Required archive path and
native backup executable. \\
Primary: \code{--database} & Required primary database path. \\
Both: \code{--cycles} & 0 means continue indefinitely; positive bounds are
useful for qualification. \\
Both: \code{--interval-ms} & 250 between cycles. \\
Both: \code{--command-timeout} & 600 seconds for native child operations. \\
Both: \code{--status-file}, \code{--ready-file} & Optional atomic status and
startup readiness files. \\
Standby: \code{--slot-root}, \code{--server-exe} & Required alternate-slot
root and native server executable. \\
Standby: \code{--listen-host} & \code{127.0.0.1}; loopback-only switch. \\
Standby: \code{--listen-port} & Required stable read endpoint port. \\
Standby: \code{--backend-port-a}, \code{--backend-port-b} & 0 selects free
backend ports. \\
Standby: \code{--maximum-clients} & 16 for backend native server capacity. \\
Standby: \code{--max-connections} & 0 means unlimited total proxy connections;
this is distinct from simultaneous backend capacity. \\
Standby: \code{--client-timeout} & 600 seconds in proxy/client handling. \\
\bottomrule
\end{longtable}

Use the atomic status files to monitor \code{healthy}, generation,
\code{sourceLsn}, \code{standbyLsn}, \code{lagBytes}, and
\code{cycleDurationMs}. Check both freshness and health; an old healthy file
from a stopped process is not a current success signal. The sidecar has a
\code{self-test} mode for its proxy/backend lifecycle. That test does not
replace a real database failover drill.

\sourcecode{tools/replication/minisql_hot_replica.py}{CLI, two-slot lifecycle, loopback switch, and status reports.}

\section{Promote a standby manually}

Promotion is a controlled change of authority. First stop or positively fence
the old primary so it cannot continue accepting writes. Stop the standby
server, verify the newest available archive state, and decide whether its
known replication lag is acceptable. Then run:

\begin{lstlisting}
minisql-backup.exe standby-promote C:\MiniSQL-Standby\slot-a
minisql-check.exe C:\MiniSQL-Standby\slot-a
minisqld.exe --serve C:\MiniSQL-Standby\slot-a 7432 32
\end{lstlisting}

Use the intended authenticated/TLS mode for a secured deployment instead of
the trusted demonstration mode above. Promotion publishes a promotion marker
and removes the standby marker. Recovery advances future transaction IDs
beyond observed archived transactions, avoiding reuse when the promoted
database begins writing.

Reconnect clients to the selected leader and verify one controlled write.
Existing transactions on a failed connection do not transfer to the promoted
database. Reconcile uncertain commit outcomes as described in the transaction
chapter. Rebuild the retired node from the new authoritative primary rather
than allowing both divergent copies to resume writing.

\warning{Asynchronous does not mean lossless}{A primary can acknowledge a
commit before a later shipping cycle archives it. If that primary is lost,
promotion may omit acknowledged changes that never reached the archive.
Measure and report lag; do not promise zero data loss or synchronous quorum
commit semantics.}

\sourcecode{src/minisql/tools/backup.ml}{promoteStandby.}
\sourcecode{docs/release/REPLICATION.md}{Supported failover scope and current limitations.}

\section{Automatic single-host failover and native write fencing}
\index{failover}\index{fencing}\index{HA controller}

The separate HA controller owns a local primary/backend lifecycle, a stable
loopback endpoint, a file witness, and automatic promotion. The witness
serializes leadership terms. A dedicated renewal thread updates the lease so
a long standby materialization does not by itself stop lease renewal.

The server is started with a lease path, epoch, node ID, and maximum clock
skew. It checks the persistent epoch and current lease before mutations and
again immediately before DML/DDL commit. Missing, malformed, expired, or
foreign records reject writes with error 9038. Reads and rollback remain
available for diagnosis.

\begin{lstlisting}[language=bash]
python3 tools/replication/minisql_ha_controller.py run \
  --primary-db /srv/minisql/db_ACTUAL_ID \
  --archive /srv/minisql-ha/archive \
  --slot-root /srv/minisql-ha/slots \
  --witness-dir /srv/minisql-ha/witness \
  --server-exe build/bin-linux/minisqld \
  --backup-exe build/bin-linux/minisql-backup \
  --proxy-port 7432 \
  --status-file /srv/minisql-ha/controller-status.json
\end{lstlisting}

Start this as the owner of the intended primary process; do not leave an
independently started writer running against the same directory. Use the
controller's documented disposable drill before adapting it to persistent
data. New proxy connections switch to the new leader after a promotion;
established sessions are not live-migrated and must reconnect if their backend
is retired or fails.

\begin{longtable}{p{.34\linewidth}p{.59\linewidth}}
\toprule
Controller run option & Default / meaning \\
\midrule
\endhead
\code{--primary-db}, \code{--archive}, \code{--slot-root},
\code{--witness-dir} & Required database and HA state directories. \\
\code{--server-exe}, \code{--backup-exe}, \code{--proxy-port} & Required
native executables and stable endpoint port. \\
\code{--backend-port-a}, \code{--backend-port-b} & 0 chooses free local
backend ports. \\
\code{--maximum-clients} & 32. \\
\code{--node-id} & 1; leader identity used in fencing records. \\
\code{--lease-ms} & 5000; lease duration. \\
\code{--clock-skew-ms} & 250; expiry allowance before a replacement leader
is promoted. \\
\code{--replication-ms} & 2000; replication-cycle interval. \\
\code{--command-timeout} & 300 seconds for child commands. \\
\code{--status-file} & Optional atomic controller status. \\
\code{--run-seconds} & 0 runs indefinitely. \\
\code{--max-failovers} & 0 leaves the failover count unbounded. \\
\bottomrule
\end{longtable}

The other controller modes are \code{self-test} for fencing record codecs
and \code{write-epoch --database path --epoch n --node-id n} for offline
epoch administration. They are implementation/qualification tools; changing
epochs without understanding the active lease and owner can revoke or confuse
leadership. Prefer the controller-owned workflow.

\subsection{What the witness does not provide}

A local file witness is not distributed consensus. The directory must have
one controller writer and reliable atomic-replacement semantics. Copying the
witness independently to multiple machines does not preserve one shared
leadership decision. A multiple-host design needs an external quorum-based
lease/leadership service and fencing at an authoritative write boundary.
The supplied controller does not add synchronous replication or zero-loss
commits.

The five-second lease and 250-millisecond skew allowance are operational
defaults, not universal safety guarantees for arbitrary clock behavior or
network filesystems. Qualify the actual host, process supervision, witness
filesystem, restart behavior, and client reconnection policy together.

\sourcecode{tools/replication/minisql_ha_controller.py}{Controller modes, leases, promotion, and stable endpoint.}
\sourcecode{src/minisql/server/database_manager.ml}{validateWriteFence and epoch/lease record checks.}
\sourcecode{docs/formats/ha-fencing-records-v1.md}{Persistent fencing-record formats.}

\section{Change page size with an offline rewrite}
\index{migration}\index{page size}

Changing \code{databaseDefaults.pageSize} in JSON does not rewrite a
database. The short migration command only checks whether the requested size
already matches:

\begin{lstlisting}
minisql-migrate.exe C:\MiniSQL-Data\db_ACTUAL_ID 8192
\end{lstlisting}

If the source already uses 8192, it reports success with no change. If it
uses another supported size, it reports that an offline rewrite is required
and does not modify the source. Request that rewrite explicitly:

\begin{lstlisting}
minisql-migrate.exe --rewrite C:\MiniSQL-Data\db_ACTUAL_ID `
  C:\MiniSQL-Migrated learning_8k 8192
\end{lstlisting}

The tool opens the stopped source, creates a new database under a staging
root, clones catalog/schema/security state, reads and rewrites rows at the
new page size, rebuilds indexes, checks the result, and publishes a new
UUID-based target path. It prints that actual path plus row/index counts.
The logical name is the supplied target name; the database identity is new.
The source remains available for rollback.

Supported sizes remain 4096, 8192, 16384, and 32768. Current copy code
materializes each source table's rows for rewriting, so plan memory as well
as additional disk capacity. It preserves the supported schema extensions,
including views, sequences, triggers, and procedure definitions, through the
cloned schema state, but it does not copy the old physical WAL history or
automatically preserve the source's TDE envelope. Enable/verify encryption
on the new target if required, refresh optimizer statistics as appropriate,
and create a fresh replication base.

This command accepts only databases the current engine can open. It cannot
convert a rejected format-1 database into format 2. For that version boundary,
use the old matching engine to extract logical data and rebuild it in a new
current-format database. Compare actual rows and schema objects before
retiring the old copy.

\sourcecode{src/minisql/tools/migrate.ml}{plan, run, copyDatabaseState, and rewrite.}
\sourcecode{src/apps/minisql_migrate/main.ml}{Public migration arguments and success output.}

\section{Complete backup, checker, and migration command reference}

\begin{lstlisting}
minisql-check --version
minisql-check --m0-self-test
minisql-check <database>

minisql-backup --version
minisql-backup --m0-self-test
minisql-backup backup <database> <new-backup>
minisql-backup backup-encrypted <database> <new-backup> <key-file>
minisql-backup restore <backup> <new-database>
minisql-backup restore-encrypted <backup> <new-database> <key-file>
minisql-backup archive-init <database> <new-archive>
minisql-backup archive-wal <stopped-database> <archive>
minisql-backup archive-wal-live <database> <archive>
minisql-backup archive-verify <archive>
minisql-backup restore-pitr <archive> <new-database> <lsn|latest>
minisql-backup standby-materialize <archive> <new-standby>
minisql-backup standby-refresh <archive> <stopped-standby>
minisql-backup standby-promote <stopped-standby>

minisql-migrate --version
minisql-migrate --m0-self-test
minisql-migrate <stopped-database> <page-size>
minisql-migrate --rewrite <stopped-source> <target-root>
  <target-name> <page-size>
\end{lstlisting}

Use \path{.exe} on Windows. The bracketed \code{lsn|latest} position means
choose a numeric LSN or the literal \code{latest}; do not pass the bar.
Wrapped reference lines are one command. These tools return 0 on success,
1 for a reported operational failure, and 2 for invalid usage. Preserve the
printed error code/message, rather than treating every nonzero exit as a
generic copy failure.

\sourcecode{src/apps/minisql_backup/main.ml}{Complete backup/archive/standby dispatch.}
\sourcecode{src/apps/minisql_check/main.ml}{Checker dispatch and warnings.}
\sourcecode{src/apps/minisql_migrate/main.ml}{Migration dispatch.}

\section{A recovery rehearsal from beginning to end}

For a disposable learning database, insert a distinctive baseline record,
stop the server, and create a base/archive. Restart it, insert a second record,
and export a live WAL prefix. Restore the base to one fresh path and the
latest archive to another. Verify that the base has only the baseline record
and the latest restore includes the later committed record. Run the checker
on each stopped copy.

Next materialize a standby, query it through a standby listener, and confirm
that a write is rejected. Stop the original primary and standby listener,
promote the standby, restart it in the intended writable mode, and commit a
third record. This teaches the distinction between a verified backup, an
applied archive generation, a read-only standby, and an authoritative writable
leader.

For automated qualification, use the repository's production failover,
automatic-fencing, and automatic-controller drills in new empty work roots.
The fencing drill deliberately leaves an old leader alive and verifies write
rejection; the controller drill terminates a controller-owned leader and
checks promotion. Keep the reports, repeat them for the deployed revision,
and measure recovery time on data volumes representative of the actual
installation.

\sourcecode{tests/ha/production_failover_drill.py}{Live shipping and manual promotion drill.}
\sourcecode{tests/ha/automatic_fencing_drill.py}{Stale-writer rejection and fenced leadership transition.}
\sourcecode{tests/ha/automatic_controller_live.py}{Controller-owned leader failure and automatic promotion.}
